\documentclass[aspectratio=169,11pt]{beamer}

\usetheme{Madrid}
\usecolortheme{default}
\usefonttheme{professionalfonts}
\setbeamertemplate{navigation symbols}{}
\setbeamertemplate{footline}[frame number]

\usepackage{amsmath, amssymb, booktabs}

\title{Disease Exposure, Environmental Health, Demographic Change, And Labor-Market Adjustment}
\subtitle{Incidence, Reallocation, Migration, Aging, And Worker Welfare}
\author{Special Topic 6: Labor Market, Health, and Population -- Week 5}
\date{}

\begin{document}

\begin{frame}
  \titlepage
\end{frame}

\begin{frame}{Central question}
\begin{itemize}
    \item How do disease, environmental exposure, and demographic structure reshape workers, firms, local labor markets, and aggregate adjustment?
    \item The labor objects are employment, hours, wages, productivity, migration, retirement, technology adoption, job quality, and welfare.
    \item Week 5 discipline: distinguish short-run disruption, medium-run reallocation, and long-run distributional or institutional effects.
\end{itemize}
\end{frame}

\begin{frame}{Disease, environment, and demography as labor shocks}
\small
\begin{tabular}{p{0.25\linewidth} p{0.34\linewidth} p{0.28\linewidth}}
\toprule
Shock & Immediate labor margin & Adjustment margin \\
\midrule
Infectious disease & absences, closures, risk, care & sectoral recovery, safety, migration \\
Pollution / heat & fatigue, productivity, hours & shifts, task redesign, sorting \\
Mortality crisis & scarcity, household disruption & wages, institutions, skill premia \\
Aging / fertility decline & participation, retirement, care & automation, migration, job design \\
\bottomrule
\end{tabular}
\vspace{0.3cm}
\begin{itemize}
    \item Incidence can appear in wages, hours, output, rents, firm exit, migration, retirement, or hidden welfare losses.
\end{itemize}
\end{frame}

\begin{frame}{A labor-market adjustment system}
\small
\[
Y_{rgt}=f(H_{rgt},E_{rt},D_{rt},F_{rt},M_{rt},P_{rt})
\]
\begin{itemize}
    \item $H$: health capacity, morbidity, mortality risk, and fatigue.
    \item $E$: disease, pollution, heat, workplace, or neighborhood exposure.
    \item $D$: age structure, fertility, dependency ratios, and composition.
    \item $F$: firm demand, production technology, scheduling, and entry or exit.
    \item $M$: migration, commuting, remote work, occupation switching, and immobility.
    \item $P$: public health, shutdowns, insurance, leave, pensions, and infrastructure.
\end{itemize}
\end{frame}

\begin{frame}{Historical health episodes and labor outcomes}
\small
\begin{tabular}{p{0.19\linewidth} p{0.27\linewidth} p{0.22\linewidth} p{0.20\linewidth}}
\toprule
Episode & Shock object & Labor margin & Persistence question \\
\midrule
Black Death & mortality and scarcity & wages, bargaining, institutions & inequality and labor rules \\
Tuberculosis & public-health capacity & cohort health, productivity, mobility & human capital and urban mortality \\
HIV/AIDS & morbidity and mortality & wages, employment, care, skill premia & household and local adjustment \\
COVID & disease plus policy shock & employment, remote work, care, sectors & reallocation and scarring \\
\bottomrule
\end{tabular}
\vspace{0.3cm}
\begin{itemize}
    \item Historical episodes teach the same labor question: who bears incidence, who adjusts, and what persists?
\end{itemize}
\end{frame}

\begin{frame}{COVID in historical perspective}
\begin{itemize}
    \item COVID is a modern epidemic labor shock with fast data, strong policy response, sectoral concentration, remote work, and school or care disruptions.
    \item Like older disease shocks, it bundles health risk, labor demand, public institutions, household care, migration, and inequality.
    \item Employment losses can reflect infection risk, shutdown policy, customer demand, firm closure, transport risk, or family constraints.
    \item High-frequency payroll, transaction, mobility, and vacancy data help with timing but do not measure worker welfare by themselves.
\end{itemize}
\end{frame}

\begin{frame}{Environmental health and hidden incidence}
\small
\begin{tabular}{p{0.24\linewidth} p{0.32\linewidth} p{0.30\linewidth}}
\toprule
Exposure & Visible outcome & Hidden labor incidence \\
\midrule
Air pollution & hours, absences, output & fatigue, errors, lower job quality \\
Heat & time allocation, productivity & unsafe work, shift timing, recovery \\
Wildfire smoke & attendance and location choice & risk, caregiving, school disruption \\
Workplace exposure & safety, turnover, claims & stress, sorting, compensating gaps \\
\bottomrule
\end{tabular}
\vspace{0.25cm}
\begin{itemize}
    \item Wages and employment can miss productivity loss, presenteeism, risk, and the lower real value of work.
\end{itemize}
\end{frame}

\begin{frame}{Health inequality, migration, and adjustment}
\begin{itemize}
    \item Exposure differs by occupation, neighborhood, housing, transport, race, wealth, age, migrant status, and local institutions.
    \item Adjustment capacity differs by remote-work feasibility, savings, insurance, credentials, legal status, family obligations, and destination access.
    \item Migration is both an outcome and a selection problem: movers and stayers reveal different constraints.
    \item Immobility is a labor-market object when workers remain in risky jobs or places because adjustment is too costly.
    \item Global and development settings make informality, crowding, weak insurance, and migrant vulnerability central.
\end{itemize}
\end{frame}

\begin{frame}{Aging and demographic change}
\small
\begin{tabular}{p{0.24\linewidth} p{0.34\linewidth} p{0.28\linewidth}}
\toprule
Demographic force & Labor-market margin & Firm/place response \\
\midrule
Aging & participation, retirement, health limits & automation, accommodation, job redesign \\
Fertility decline & future worker supply and care burden & recruitment, human capital, migration \\
Dependency structure & pensions, fiscal pressure, care demand & elder care, benefits, late-career work \\
Composition change & skill and age mix & wage structure, sectoral change \\
\bottomrule
\end{tabular}
\vspace{0.3cm}
\begin{itemize}
    \item Demography affects labor markets through supply, demand, care, technology, migration, and retirement margins at once.
\end{itemize}
\end{frame}

\begin{frame}{Methods and data}
\begin{itemize}
    \item Cohort-by-place designs: historical epidemics, fetal exposure, public-health campaigns, long-run outcomes.
    \item Event studies: outbreaks, shutdowns, heat waves, pollution spikes, mortality shocks, firm or regional exposure.
    \item Exposure gradients: monitors, weather, satellite data, reporting laws, regulations, occupation risk, and policy timing.
    \item Panels and linked data: earnings, payroll, UI, claims, mortality, attendance, output, firms, migration, and commuting.
    \item Threats: survival selection, avoidance, endogenous residence, migration selection, local trends, bundled policies, and spillovers.
\end{itemize}
\end{frame}

\begin{frame}{Research Lab design}
\begin{itemize}
    \item Primary anchor: COVID labor-market incidence and unequal employment effects by worker type and occupation.
    \item Challenge anchor: HIV/AIDS mortality, development, skill premia, and longer-run distributional adjustment.
    \item Reproduce: build a synthetic COVID-style incidence profile across occupation, exposure, care burden, migrant status, and demography.
    \item Diagnose: classify disease risk, policy and demand collapse, caregiving, sorting, migration, firm response, and hidden welfare incidence.
    \item Transfer: redesign the workflow for plague, tuberculosis, HIV/AIDS, pollution, heat, migrant vulnerability, or aging and automation.
\end{itemize}
\end{frame}

\begin{frame}{Research design checklist}
\begin{enumerate}
    \item Name the shock: disease, environment, mortality, aging, fertility, demography, or policy.
    \item Name the margin: work capacity, employment, hours, productivity, wages, migration, retirement, automation, or welfare.
    \item State the timing: short-run disruption, medium-run reallocation, or long-run persistence.
    \item Name the unit: worker, household, firm, occupation, region, cohort, or sector.
    \item Say what remains bundled: policy, demand, avoidance, migration, survival, or general-equilibrium response.
\end{enumerate}
\end{frame}

\begin{frame}{Bridge to Week 6}
\begin{itemize}
    \item Week 5 supplies the scaled-up objects: disease, environment, demography, migration, firms, places, and welfare.
    \item Week 6 turns the full course into research designs with explicit labor outcomes, counterfactuals, mechanisms, and welfare objects.
    \item The through-line is disciplined incidence: who is exposed, who adjusts, who is constrained, and what wages or employment miss.
\end{itemize}
\end{frame}

\begin{frame}{Takeaways}
\begin{itemize}
    \item Disease, environmental exposure, and demographic change are labor-market shocks because they alter capacity, demand, productivity, mobility, and welfare.
    \item Historical epidemics and COVID belong in one incidence literature when we compare timing, adjustment, institutions, and inequality.
    \item Health inequality is central: unequal exposure plus unequal adjustment capacity determines who bears labor-market costs.
    \item Strong designs name the shock, margin, timing, unit, data, threat, and welfare object.
\end{itemize}
\end{frame}

\end{document}
